\documentclass[12pt,a4paper]{article}
\usepackage{ugmath}
\usepackage{float}
\usepackage{placeins}
\usepackage{array}
\usepackage[labelfont=bf,labelsep=space]{caption}
\newcommand{\paren}[1]{\left( #1 \right)}
\newcommand{\alumno}{Ricardo León Martínez}
\newcommand{\materia}{Elementos de Probabilidad y Estadistica}
\newcommand{\profesor}{Claudia Reynoso Alcántara}
\newcommand{\tarea}{Tarea 5}
\newcommand{\fecha}{18/3/2026}

\begin{document}

    \begin{center}
        {\large\textbf{UNIVERSIDAD DE GUANAJUATO}}\\[0.3cm]
        {\normalsize\textbf{DIVISIÓN DE CIENCIAS NATURALES Y EXACTAS}}\\
        {\normalsize\textbf{CAMPUS GUANAJUATO}}\\[1cm]

        {\Large\textbf{\tarea\ (\materia)}}\\[1cm]
    \end{center}

    \noindent
    \textbf{Nombre:} \alumno 
    \hfill 
    \textbf{Fecha:} \fecha 
    \hfill 
    \textbf{Calificación:} \rule{3cm}{0.4pt}

    \vspace{0.3cm}

    \begin{mdframed}[style=mdbluebox,frametitle={Ejercicio 1}]
        En cierta encuesta participaron 1000 personas a las que se les preguntó sobre su ocupación,
        estado civil y nivel educativo. De las personas encuestadas y los datos reportados se sabe que 312
        son profesionistas, 470 están casadas, 525 poseen títulos universitarios, 42 son profesionistas y poseen
        títulos universitarios, 147 están casadas y poseen títulos universitarios, 86 son profesionistas y están
        casadas y 25 son profesionistas casadas que poseen títulos universitarios. Demuestre que los números
        reportados deben ser incorrectos.
    \end{mdframed}

    \begin{proof}
        Sea $U$ el universo de las 1000 personas encuestadas, y definimos los subconjuntos
        \begin{align*}
            P&:=\{\text{profesionistas}\}\\
            C&:=\{\text{personas casadas}\}\\
            T&:=\{\text{personas con título universitario}\}.
        \end{align*}
        Se nos da lo siguiente
        \begin{align*}
            \abs{P}&=312,\quad\abs{C}=470,\quad\abs{T}=525,\\
            \abs{P\cap T}&=42,\quad\abs{C\cap T}=147,\\
            \abs{P\cap C}&=86,\quad\abs{P\cap C\cap T}=25.
        \end{align*}
        Por el principio de inclusión-exclusión se tiene que
        \begin{equation*}
            \abs{P\cup C\cup T}=\abs{P}+\abs{C}+\abs{T}-\abs{P\cap C}-\abs{P\cap T}-\abs{C\cap T}+\abs{P\cap C\cap T}
        \end{equation*}
        Sustituyendo
        \begin{align*}
            \abs{P\cup C\cup T}&=312+470+525-86-42-147+25.\\
            &=1057
        \end{align*}
        Por definición, $P\cup C\cup T\subseteq U$, por lo que necesariamente
        \begin{equation*}
            \abs{P\cup C\cup T}\leq\abs{U}=1000.
        \end{equation*}
        Pero, hemos obtenido
        \begin{equation*}
            \abs{P\cup C\cup T}=1057\gt1000,
        \end{equation*}
        lo cual es imposible.
    \end{proof}

    \begin{mdframed}[style=mdbluebox,frametitle={Ejercicio 2}]
        Un examen contiene 11 preguntas en total a contestar con falso o verdadero. El examen
        resulta aprobado si se responden correctamente al menos 6 de las 11 preguntas. Si contestamos las
        preguntas al azar, ¿cuál es la probabilidad de aprobar el examen?
    \end{mdframed}

    \begin{solucion}
        Sea $X$ la variable que cuenta el número de respuestas correctas. Aprobar el
        examen equivale a que $X\geq6$. Por lo tanto
        \begin{equation*}
            \mathbb{P}(X\geq6)=\sum_{k=6}^{11}\binom{11}{k}\paren{\frac{1}{2}}^{11}
        \end{equation*}
        Calculamos los coeficientes binomiales
        \begin{align*}
            \binom{11}{6}&=462,\\
            \binom{11}{7}&=330,\\
            \binom{11}{8}&=165,\\
            \binom{11}{9}&=55,\\
            \binom{11}{10}&=11,\\
            \binom{11}{11}&=1.
        \end{align*}
        Sumando
        \begin{equation*}
            462+330+165+55+11+1=1024.
        \end{equation*}
        Así,
        \begin{equation*}
            \mathbb{P}(X\geq6)=\frac{1024}{2^{11}}=\frac{1024}{2048}=\frac{1}{2}.
        \end{equation*}
    \end{solucion}

    \begin{mdframed}[style=mdbluebox,frametitle={Ejercicio 3}]
        Describa espacios muestrales para cada uno de los siguientes casos.
        \begin{itemize}
            \item[(a)] Calificaciones de una clase con 37 estudiantes en un examen.
            \item[(b)] Medición de la temperatura a medio día en una estación meteorológica.
            \item[(c)] Medición de las velocidades de automóviles pasando por un punto dado.
            \item[(d)] Lanzamientos infinitos de una moneda.
        \end{itemize}
    \end{mdframed}

    \begin{solucion}
        \begin{enumerate}
            \item[(a)] Si las calificaciones pertenecen, por ejemplo, a un subconjunto $A\subseteq\mathbb{R}$
            (por ejemplo $A=[0,10]$ o bien un conjunto discreto como $\{0,1,\dots,10\}$), entonces
            cada resultado del experimento consiste en una 37-tupla de calificaciones
            \begin{equation*}
                \Omega=A^{37}=\{(x_{1},x_{2},\dots,x_{37}):x_{i}\in A\}.
            \end{equation*}
            \item[(b)] Suponiendo que la estacion mide a la misma hora, con el mismo instrumento de medición,
            y la misma precision. Si el termometro mide en decimas de grado dentro de un rango fisico $[a,b]$,
            entonces el espacio muestral es finito o numerable
            \begin{equation*}
                \Omega=\{a,a+0.1,a+0.2,\dots,b\}.
            \end{equation*}
            \item[(c)] Suponiendo las mismas cosas que el anterior, entonces, el espacio muestral queda
            \begin{equation*}
                \Omega=\{1,2,\dots,V_{\text{max}}\}.
            \end{equation*}
            \item[(d)] Cada lanzamiento produce cara (C) o cruz (X). Un resultado del experimento es una
            sucesión infinita. Por lo tanto,
            \begin{equation*}
                \Omega=\{(x_{1},x_{2},x_{3},\dots):x_{i}\in\{C,X\}\}.
            \end{equation*}
        \end{enumerate}
    \end{solucion}

    \begin{mdframed}[style=mdbluebox,frametitle={Ejercicio 4}]
        Fulgencio y Yesrael acuerdan realizar una competencia de esgrima con duelos consecutivos
        en la que el ganador de la competencia es el primero en salir victorioso en dos duelos consecutivos.
        En cada duelo, la probabilidad de que Fulgencio salga victorioso es $p$ con $0<p<1$, mientras que la
        probabilidad de victoria de Yesrael es $q = 1-p$.
        \begin{itemize}
            \item[(a)] Describa un espacio muestral para este proceso.
            \item[(b)] Calcule la probabilidad de que la competencia termine al cabo de $k$ duelos con $k=2,3,\dots$.
        \end{itemize}
    \end{mdframed}

    \begin{solucion}
        \begin{enumerate}
            \item[(a)] El experimento consiste en observar la sucesión de resultados hasta que aparece por primera vez
            dos victorias consecutivas de alguno de los jugadores. Por tanto el espacio muestra esta formado
            por todas las sucesiones finitas de simbolos $\{F, Y\}$ que terminan en dos resultados iguales
            consecutivos, pero que antes de ese punto no contienen dos iguales consecutivos.
            \item[(b)] Para que el juego termine exactamente en el duelo $k$:
            \begin{itemize}
                \item Los primero $k-1$ resultados deben alternar estrictamente.
                \item En el paso $k$, se repite el resultado anterior.
            \end{itemize}
            Esto implica que
            \begin{itemize}
                \item Existe exactamente una sucesion que termina en $F$.
                \item Existe exactamente una sucesion que termina en $Y$.
            \end{itemize}
            Si $k=2m$, entonces las dos sucesiones posibles son
            \begin{equation*}
                Y,F,Y,F,\dots,F,F\quad\text{ y }\quad F,Y,F,Y,\dots,Y,Y.
            \end{equation*}
            Cada una tiene $m$ apariciones de $F$ y $m$ de $Y$. Por la independencia
            \begin{equation*}
                \mathbb{P}=p^{m}q^{m}.
            \end{equation*}
            Como hay dos sucesiones
            \begin{equation*}
                \mathbb{P}(K=2m)=2(pq)^{m}
            \end{equation*}
            Si $k=2m+1$, las sucesiones son
            \begin{equation*}
                F,Y,F,Y,\dots,F,F\quad\text{ y }\quad Y,F,Y,F,\dots,Y,Y.
            \end{equation*}
            En la primera hay $m+1$ $F$ y $m$ $Y$, en la segunda es al reves. Por lo tanto
            \begin{equation*}
                \mathbb{P}(K=2m+1)=p^{m+1}q^{m}+q^{m+1}p^{m}
            \end{equation*}
        \end{enumerate}
    \end{solucion}

    \begin{mdframed}[style=mdbluebox,frametitle={Ejercicio 5}]
        Tres jugadores, a los cuales identificamos por Jugador 1, Jugador 2 y Jugador 3, toman
        turnos consecutivos para lanzar una moneda con caras sol y águila. La moneda podría estar desbalanceada,
        por lo que debemos asumir que la probabilidad de que salga águila es $p$ con $0<p<1$, mientras
        que la probabilidad de sol $q = 1-p$. El primer lanzamiento lo realiza el Jugador 1, el segundo el Jugador
        2, el tercero el Jugador 3, el cuarto el Jugador 1, el quinto el Jugador 2, y así sucesivamente. Gana el
        primer jugador que obtenga un sol.
        \begin{itemize}
            \item[(a)] Describa un espacio muestral para este proceso.
            \item[(b)] Describa el evento de que el Jugador 1 gane y calcule su probabilidad.
            \item[(c)] Describa el evento de que el Jugador 2 gane y calcule su probabilidad.
            \item[(d)] Describa el evento de que ni el Jugador 1 ni tampoco el Jugador 2 gana y calcule su probabilidad.
        \end{itemize}
    \end{mdframed}

    \begin{solucion}
        \begin{enumerate}
            \item[(a)] El experimento consiste en lanzar la moneda hasta la primera aparición de $S$. Por tanto,
            \begin{equation*}
                \Omega=\{S,AS,AAS,AAAS,\dots\}.
            \end{equation*}
            Cada $\omega\in\Omega$ es una sucesión finita que termina en el primer $S$.
            \item[(b)] El jugador 1 lanza en los tiempos $1,4,7,\dots$, es decir en los indices
            $k\equiv1\pmod{3}$. Entonces,
            \begin{equation*}
                E_{1}=\{\omega\in\Omega:\abs{\omega}=3m+1,m\geq0\}.
            \end{equation*}
            Para $k=3m+1$, la probabilidad de la trayectoria correspondiente es
            \begin{equation*}
                \mathbb{P}(\Omega_{3m+1})=p^{3m}q.
            \end{equation*}
            Por tanto,
            \begin{equation*}
                \mathbb{P}(E_{1})=\sum_{m=0}^{\infty}p^{3m}q=q\sum_{m=0}^{\infty}(p^{3})^{m}.
            \end{equation*}
            Como $0\lt p\lt1$, la serie geometrica converge y
            \begin{equation*}
                \mathbb{P}(E_{1})=\frac{q}{1-p^{3}}.
            \end{equation*}
            \item[(c)] El jugador 2 lanza en los tiempos $2,5,8,\dots$, es decir $k\equiv2\pmod{3}$.
            \begin{equation*}
                E_{2}=\{\omega\in\Omega:\abs{\omega}=3m+2m,m\geq0\}
            \end{equation*}
            Para $k=3m+2$, la probabilidad es
            \begin{equation*}
                \mathbb{P}(\Omega_{3m+2})=p^{3m+1}q.
            \end{equation*}
            Así,
            \begin{equation*}
                \mathbb{P}(E_{2})=\sum_{m=0}^{\infty}p^{3m+1}q=pq\sum_{m=0}^{\infty}=\frac{pq}{1-p^{3}}.
            \end{equation*}
            \item[(d)] Este evento corresponde a que gana el jugador 3:
            \begin{equation*}
                E_{3}=\{\omega\in\Omega:\abs{\omega}=3m+3,m\geq0\}.
            \end{equation*}
            Para $k=3m+3$,
            \begin{equation*}
                \mathbb{P}(\Omega_{3m+3})=p^{3m+2}q.
            \end{equation*}
            Por tanto,
            \begin{equation*}
                \mathbb{P}(E_{3})=\sum_{m=0}^{\infty}p^{3m+2}q=p^{2}q\sum_{m=0}^{\infty}(p^{3})^{m}=\frac{p^{2}q}{1-p^{3}}.
            \end{equation*}
        \end{enumerate}
    \end{solucion}

    \begin{mdframed}[style=mdbluebox,frametitle={Ejercicio 6}]
        Sea $\Omega$ un conjunto no vacío. Recordemos que $\mathcal{F}\subset 2^{\Omega}$ es una
        $\sigma$-álgebra (de subconjuntos) de $\Omega$ si satisface que:
        \begin{itemize}
            \item[(I)] $\Omega \in \mathcal{F}$,
            \item[(II)] $A\in\mathcal{F}$ implica $A^{c}\in\mathcal{F}$,
            \item[(III)] $\{A_k\}_{k=1}^{\infty}\subset \mathcal{F}$ implica $\displaystyle\bigcup_{k=1}^{\infty}A_k \in \mathcal{F}$.
        \end{itemize}
        Más aún, dada una colección $\{S_i\}_{i\in I}\subset 2^{\Omega}$ de subconjuntos de $\Omega$, la $\sigma$-álgebra
        generada por $\{S_i\}_{i\in I}$ es la única $\sigma$-álgebra $\mathcal{F}$ que cumple que:
        \begin{itemize}
            \item[(a)] $\mathcal{F}$ es una $\sigma$-álgebra, es decir, satisface (I), (II) y (III);
            \item[(b)] $S_i\in\mathcal{F}$ para todo $i\in I$;
            \item[(c)] si $\mathcal{F}'$ es una $\sigma$-álgebra tal que $S_i\in\mathcal{F}'$ para todo $i\in I$, entonces $\mathcal{F}\subset \mathcal{F}'$.
        \end{itemize}

        Demuestre la existencia de la $\sigma$-álgebra generada por una colección arbitraria $\{S_i\}_{i\in I}\subset 2^{\Omega}$ de subconjuntos de $\Omega$. Para ello, siga los siguientes pasos:
        \begin{itemize}
        \item[(i)] Note que $2^{\Omega}$ es una $\sigma$-álgebra tal que $S_i\in 2^{\Omega}$ para todo $i\in I$. De este modo, el conjunto
        \[
            \Gamma=\{\mathcal{F}'\subset 2^{\Omega}:\mathcal{F}'\text{ es una }\sigma\text{-álgebra tal que }S_i\in\mathcal{F}'\text{ para todo }i\in I\}
        \]
        es no vacío.
        \item[(ii)] La intersección de todos los elementos de $\Gamma$ pertenece a $\Gamma$, es decir,
        \[
            \bigcap_{\mathcal{F}'\in\Gamma}\mathcal{F}'.
        \]
        \item[(iii)] La $\sigma$-álgebra generada por $\{S_i\}_{i\in I}$ debe ser $\displaystyle\bigcap_{\mathcal{F}'\in\Gamma}\mathcal{F}'$.
        \end{itemize}
        \textbf{Observación.} De manera general, si $\{\mathcal{F}'_j\}_{j\in J}$ es una familia de $\sigma$-álgebras de $\Omega$, entonces
        \[
            \bigcap_{j\in J}\mathcal{F}'_j
        \]
        también es una $\sigma$-álgebra de $\Omega$.
    \end{mdframed}

    \begin{mdframed}[style=mdbluebox,frametitle={Ejercicio 7}]
        Sea $(\Omega,\mathcal{F},P)$ un espacio de probabilidad. Supongamos que
        $\{A_n\}_{n=1}^{\infty}$ y $\{B_n\}_{n=1}^{\infty}$ son dos colecciones numerables de eventos, es decir,
        $A_n,B_n\in\mathcal{F}$ para $n=1,2,\dots$. Demuestre que
        \[
            \lim_{n\to\infty}P(B_n)=1 \quad \text{implica} \quad
            \lim_{n\to\infty}\big(P(A_n)-P(A_n\cap B_n)\big)=0.
        \]
        \textbf{Observación.} Es posible que $\lim_{n\to\infty}P(A_n)$ no exista.
    \end{mdframed}

    \begin{proof}
        Observemos que para todo $n\in\mathbb{N}$,
        \begin{equation*}
            P(A_{n})-P(A_{n}\cap B_{n})=P(A_{n}\cap B_{n}^{c}),
        \end{equation*}
        ya que $(A_{n}\cap B_{n})\cup(A_{n}\cap B_{n}^{c})$ es union disjunta. Por lo tanto,
        basta demostrar que
        \begin{equation*}
            \lim_{n\to\infty}P(A_{n}\cap B_{n}^{c})=0.
        \end{equation*}
        Para todo $n$, se tiene la inclusión
        \begin{equation*}
            A_{n}\cap B_{n}^{c}\subseteq B_{n}^{c},
        \end{equation*}
        y por monotonia de la probabilidad
        \begin{equation*}
            0\leq P(A_{n}\cap B_{n}^{c})\leq P(B_{n}^{c}).
        \end{equation*}
        Como por hipotesis
        \begin{equation*}
            \lim_{n\to\infty}P(B_{n})=1,
        \end{equation*}
        se sigue que
        \begin{equation*}
            \lim_{n\to\infty}P(B_{n}^{c})=\lim_{n\to\infty}(1-P(B_{n}))=0.
        \end{equation*}
        Entonces, por el teorema del sandwich
        \begin{equation*}
            0\leq P(A_{n}\cap B_{n}^{c})\leq P(B_{n}^{c})\to0
        \end{equation*}
        Lo cual implica
        \begin{equation*}
            \lim_{n\to\infty}P(A_{n}\cap B_{n}^{c})=0.
        \end{equation*}
        Así,
        \begin{equation*}
            \lim_{n\to\infty}(P(A_{n})-P(A_{n}\cap B_{n}))=0.
        \end{equation*}
    \end{proof}

    \begin{mdframed}[style=mdbluebox,frametitle={Ejercicio 8}]
        Sea $(\Omega,\mathcal{F},P)$ un espacio de probabilidad. Pruebe que se cumple la siguiente desigualdad de Bonferroni:
        \[
        P\left(\bigcup_{k=1}^{n}A_k\right)
        \ge
        \sum_{k=1}^{n}P(A_k)
        -
        \sum_{1\le k<\ell\le n}P(A_k\cap A_\ell).
        \]
        \textbf{Observación/Sugerencia.} Lo anterior es una versión truncada del principio de inclusión-exclusión,
        el cual se prueba por inducción. Recordemos además que para cualesquiera eventos
        $B_1,B_2,\dots,B_n\in\mathcal{F}$ se satisface que
        \[
        P(B_1\cup B_2\cup \cdots \cup B_n)
        \le
        P(B_1)+P(B_2)+\cdots+P(B_n).
        \]
    \end{mdframed}

    \begin{proof}
        Procederemos por inducción en $n$. Caso base $n=2$. Para dos eventos se tiene
        la igualdad de inclusión-exclusión
        \begin{equation*}
            P(A_{1}\cup A_{2})=P(A_{1})+P(A_{2})-P(A_{1}\cap A_{2})
        \end{equation*}
        Por lo que la desigualdad es valida. Supongamos que para algun $n\in\mathbb{N}$
        se cumple
        \begin{equation*}
            P\paren{\bigcup_{k=1}^{n}A_{k}}\geq\sum_{k=1}^{n}P(A_{k})-\sum_{1\leq k\leq\ell\leq n}P(A_{k}\cap A_{\ell}).
        \end{equation*}
        Definamos
        \begin{equation*}
            B:=\bigcup_{k=1}^{n}A_{k}.
        \end{equation*}
        Entonces,
        \begin{equation*}
            P\paren{\bigcup_{k=1}^{n+1}A_{k}}=P(B\cup A_{n+1})=P(B)+P(A_{n+1})-P(B\cap A_{n+1}).
        \end{equation*}
        Observemos que
        \begin{equation*}
            B\cap A_{n+1}=\bigcup_{k=1}^{n}(A_{k}\cap A_{n+1}).
        \end{equation*}
        Por aditividad
        \begin{equation*}
            P(B\cap A_{n+1})\leq\sum_{k=1}^{n}P(A_{k}\cap A_{n+1}).
        \end{equation*}
        Sustituyendo en la expresion anterior 
        \begin{equation*}
            P\paren{\bigcup_{k=1}^{n+1}A_{k}}\geq P(B)+P(A_{n+1})-\sum_{k=1}^{n}P(A_{k}\cap A_{n+1}).
        \end{equation*}
        Aplicando ahora la hipotesis de induccion a $P(B)$
        \begin{equation*}
            P(B)\geq\sum_{k=1}^{n}P(A_{k})-\sum_{1\leq k\lt\ell\leq n}P(A_{k}\cap A_{\ell}).
        \end{equation*}
        Por lo tanto
        \begin{equation*}
            P\paren{\bigcup_{k=1}^{n}A_{k}}\geq\sum_{k=1}^{n}P(A_{k})-\sum_{1\leq k\lt\ell\leq n}P(A_{k}\cap A_{\ell})
            +P(A_{n+1})-\sum_{k=1}^{n}P(A_{k}\cap A_{n+1}).
        \end{equation*}
        Reordenando
        \begin{equation*}
            P\paren{\bigcup_{k=1}^{n+1}A_{k}}\geq\sum_{k=1}^{n+1}P(A_{k})-\sum_{1\leq k\lt\ell\leq1}P(A_{k}\cap A_{\ell}).
        \end{equation*}
    \end{proof}
\end{document}